\chapter{Discussie}
\label{ch:discussie}

\noindent In dit hoofdstuk worden de belangrijkste beperkingen en kanttekeningen
van het onderzoek gebundeld. Deze stonden eerder verspreid bij de afzonderlijke
deelonderzoeken en zijn hier samengebracht zodat hun invloed op de eindconclusie
in samenhang beoordeeld kan worden.

% ============================================================
\section{Beperkingen weerstands- en aandrijflijnberekening}
\label{subsec:holtrop-validatie}

De Holtrop-Mennen methode is een statistische regressiemethode die
geschikt is voor vroege ontwerpfasen \cite{Holtrop1984}. De methode
geeft geen garantie op een specifieke nauwkeurigheid, want dat hangt
af van hoe goed het schip lijkt op de dataset waarop de regressie
gebaseerd is. Voor \emph{'t~Jansje} zijn drie kanttekeningen van belang:

\begin{itemize}
    \item \textbf{Aanhangsels:} $R_{APP}$ is gebaseerd op een schatting van het
          natte roeroppervlak achter een skeg. Bij definitieve vaststelling van
          het aandrijfconcept dient dit nader te worden bepaald.

    \item \textbf{Ondiepwater-effect:} De methode is ontwikkeld voor diepwater.
          Het vaargebied (Nieuwe Maas, Maassluis) heeft een waterdiepte van
          5--10~m, wat bij een diepgang van 1{,}50~m een verhouding
          $h/T \approx 3$--$7$ oplevert. Bij $h/T < 4$ treden
          ondiepwatereffecten op.

    \item \textbf{Golven en wind:} De berekening geldt voor kalm water. Marges
          voor golfweerstand en windweerstand zijn toegevoegd in de
          20\%-marge uit Sectie~\ref{subsec:aandrijflijn-berekening}.
\end{itemize}

% ============================================================
\section{Beperkingen stabiliteitsberekening}
\label{sec:stab_beperkingen}

\begin{itemize}
    \item \textbf{$KG$-schatting}: Er is geen hellingsbroef gedaan met Jansje,
          waardoor $KG$ niet bekend is. Uit
          Tabel~\ref{tab:KG_gevoeligheid} blijkt echter dat $GM$ pas onder
          het IMO minimum zakt bij $KG > 2{,}04$~m. Omdat het dek op
          2{,}05~m zit, is dit onmogelijk. Onze conclusie over de stabiliteit klopt dus ,ongeacht de werkelijke waarde van $KG$.

    \item \textbf{Vrijboordmarge}: De minimale marge van 0{,}20~m is
          conservatief aangenomen. De exacte eis volgt uit de ES-TRIN/CCR
          regelgeving.

    \item \textbf{Grote hellingshoeken}: $GM$ beschrijft uitsluitend de
         stabiliteit voor kleine hellingshoeken ($\phi < 10^\circ$). Voor een volledig
          stabiliteitsoordeel is een $GZ$-diagram vereist
          \cite[p.~134]{Barras2006}.
\end{itemize}

% ============================================================
% PLACEHOLDER: voeg hier eventueel overige discussiepunten toe
% (bv. beperkingen energieopslagkeuze, aannames model, experiment).
